\chapter*{初学者术语表与链路索引}
\addcontentsline{toc}{chapter}{初学者术语表与链路索引}

\begin{keyidea}
术语表用于回查，不代替正文中的因果解释。同一个词在硬件、电路、处理器架构和
操作系统中可能处在不同层；回查时除“它叫什么”之外，还应确认它由谁建立、
保存什么、怎样访问、何时失效。
\end{keyidea}

\section*{一、电与电路}

\begin{description}[style=nextline,leftmargin=2.8em,labelindent=0em]
  \item[电压 Voltage]
  两个位置之间的电势差，单位是伏特 V。数字信号用相对参考地的电压范围表示
  逻辑状态，电源轨则用规定电压向负载提供能量；任何电压都必须说明参考点和
  器件允许范围。

  \item[电流 Current]
  单位时间通过导体截面的电荷量，单位是安培 A。电流必须形成从电源、负载到
  返回端的闭合回路，所以地和返回路径是电路的一部分，不是原理图上可忽略的
  公共装饰。

  \item[电阻 Resistance]
  对电流的阻碍，单位是欧姆 \(\Omega\)。电阻可用于限流、分压、上拉、下拉、
  反馈和信号阻尼；同样阻值放在不同位置可能承担完全不同职责。

  \item[功率 Power]
  能量转换速率，单位是瓦特 W。直流近似中 \(P=VI\)；器件和铜线的损耗转成
  热量，因此电流能力必须同时检查压降、封装、铜宽、散热和环境温度。

  \item[GND / Ground]
  电路共同电压参考与主要电流返回网络。数字地、机壳地和建筑保护地并不天然
  相同；它们怎样连接会影响安全、ESD、EMC 和测量方法。

  \item[电源轨 Power Rail]
  向一组负载提供规定电压的网络，例如 \code{VDC}、\code{+5V} 和
  \code{+3V3}。同一机器可有多条电源轨，每条都有自己的建立时机、额定负载、
  纹波和掉电条件。

  \item[逻辑电平 Logic Level]
  接收器把一段电压范围解释为 0 或 1 的规则，不是“0 V/3.3 V 两个绝对点”。
  输入低/高阈值、输出驱动能力和不可保证的中间区必须同时满足；3.3 V 与
  1.8 V 器件不能仅因都叫数字信号就直接相连。

  \item[上拉 / 下拉 Pull-up / Pull-down]
  用电阻把无人主动驱动的信号偏置到默认高/低电平。它们防止输入漂浮，也参与
  开漏总线和低有效按键；阻值要同时满足静态电流、上升时间、抗干扰和器件规格。

  \item[开漏 Open-drain]
  输出只能主动拉低、不能主动推高，释放后依靠上拉电阻回到高电平。I2C 借此
  允许多个器件共享线路并执行仲裁；没有上拉时，“释放”不会自动形成可靠高电平。

  \item[电源时序 Power Sequencing]
  多条电源、时钟、复位和 enable 按规定先后建立/撤销的过程。各电压都测得正确
  仍不表示时序正确；设备可能要求参考时钟稳定、power-good 有效后才能释放复位。

  \item[Power-good / Reset]
  Power-good 表示某电源已进入允许范围，reset 把数字状态机保持或带回规定初态。
  reset 释放只是允许器件开始初始化，不表示固件、DRAM、总线和操作系统均已就绪。

  \item[稳压器 Regulator]
  把输入转换为受控输出电压的电路。线性稳压器通过耗散多余能量调节，开关
  稳压器通过高速开关和储能元件转换；CH224K 这种 PD 协商芯片不是稳压器。

  \item[降压转换器 Buck Converter]
  用开关、电感、电容和反馈把较高直流电压转换为较低电压。原理图连通只是第一
  步，高电流热回路、SW 节点、反馈取样和器件布局直接决定输出是否稳定。

  \item[保险丝 Fuse]
  过流持续时断开或显著限制电流的保护元件。它主要限制故障能量，不能代替过压
  保护、反接保护或正确的软件电源时序。

  \item[二极管 Diode]
  主要允许电流朝一个方向流动的器件。肖特基二极管正向压降相对较低，适合简单
  电源 OR；TVS 二极管平时不参与，瞬态过压时快速导通钳位。

  \item[MOSFET]
  由栅极电压控制导通的晶体管，可用于高速开关、负载开关和反接保护。栅极、
  源极、漏极方向及体二极管都会影响断电和反灌路径，不能只把它当理想开关。

  \item[电容 Capacitor]
  储存电场能量的器件。它可做局部去耦、输出储能、反馈补偿或高速 AC 耦合；
  容量、耐压、介质、等效串联电阻和布局共同决定实际作用。

  \item[电感 Inductor]
  储存磁场能量、阻碍电流快速变化的器件。降压转换器用电感把开关脉冲平滑为
  连续输出电流，选型要考虑电感值、饱和电流、直流电阻和温升。

  \item[去耦 Decoupling]
  用靠近负载电源脚的电容提供瞬时电流并给高频噪声短回路。它不是“电源网上
  随便多放电容”；容值组合、回路面积和与参考平面的连接都重要。

  \item[ESD]
  Electrostatic Discharge，静电放电。外部人体或线缆可能带来极快高压脉冲，
  保护器件应靠近连接器并具有合适回流；高速接口还需限制保护器件寄生电容。

  \item[USB-C CC / Power Delivery]
  CC1/CC2 用于检测插头方向、角色和基础供电能力；USB Power Delivery 可在
  CC 上协商更高电压/电流。PD 协商器请求电源档位，不等于把输入电压自行稳压
  成目标电压，VBUS 主功率也不从 CC 线流过。
\end{description}

\section*{二、原理图与 PCB}

\begin{description}[style=nextline,leftmargin=2.8em,labelindent=0em]
  \item[原理图 Schematic]
  用符号和网络表达器件应该怎样电气连接。它能描述位号、引脚、数值和逻辑网络，
  但不能单独证明实际走线阻抗、机械干涉或装配质量。

  \item[PCB]
  Printed Circuit Board，印制电路板。它把原理图 net 转化为具体焊盘、铜线、
  铜面、过孔和叠层；同一原理图采用不同布局，电源、热和高速性能可能完全不同。

  \item[位号 Reference Designator]
  某块板上器件的唯一名字，例如 R43、U14、J17。首字母通常提示器件类别，
  数字用于定位；位号不是阻值、型号或引脚号。

  \item[引脚 Pin]
  器件对外的电气连接点。原理图符号 pin、封装 pad 和实物触点必须建立正确
  映射；编号必须查资料，不能凭位置或大小猜功能。

  \item[网络 Net]
  被 EDA 声明为电气相连的一组引脚。跨页同名标签仍属于同一个 net；PCB
  工具会要求这些焊盘最终由铜连通。

  \item[走线 Trace]
  PCB 铜层上实现某个网络的实际几何路径。线宽、长度、间距、换层和参考平面
  影响电阻、电感、阻抗、温升、串扰与辐射。

  \item[过孔 Via]
  把不同 PCB 铜层连接起来的金属化孔。信号换层时其回流也需邻近参考过孔；
  高频下过孔不是没有长度和寄生参数的理想节点。

  \item[焊盘 Pad]
  PCB 上供器件引脚焊接或连接的裸露铜区域。焊盘尺寸既要匹配封装，又要满足
  阻焊开窗、钢网印刷、焊接工艺和返修空间；原理图 pin 最终必须映射到正确
  pad 编号，编号颠倒会把一张逻辑正确的图制造成电气错误的板。

  \item[焊点 Solder Joint]
  焊料把器件端子与焊盘形成的机械和电气连接。焊点可能出现虚焊、桥连、空洞、
  锡量不足或热循环裂纹，所以“网络在 PCB 文件里连通”不等于实物一定导通；
  上电前仍需目检、显微检查、通断测试或自动光学/X 射线检查。

  \item[Footprint]
  器件在 PCB 上的焊盘、外形、装配与间距定义。符号正确但 footprint 引脚映射
  错误，仍会制造出电气接反的板。

  \item[铜层与叠层 Layer / Stackup]
  PCB 由信号层、电源/地平面、介质和阻焊等层叠压而成。叠层规定层数、铜厚、
  介质厚度和材料参数，决定机械厚度、回流路径及可实现阻抗；“四层板”只说明
  层数，不能代替具体的层序和板厂叠层表。

  \item[铜皮 / 电源平面 Copper Pour / Plane]
  大面积铜可承载电源或地、降低阻抗并帮助散热。它不是越大越好：孤岛、细颈、
  跨分割回流、热焊盘连接和大电流入口都会改变实际效果，必须用 DRC 和电流/
  温升分析核对，而不能只看填充后的颜色。

  \item[参考平面 Reference Plane]
  高速信号电磁场和回流所依赖的邻近连续铜面，常为 GND。信号跨越平面开槽或
  换到没有返回路径的层，会显著增大回路和干扰。

  \item[阻抗 Impedance]
  交流或高速情况下电压与电流的综合关系，包含电阻、电感和电容效应。传输线
  特性阻抗由 PCB 几何和材料决定，不能用普通万用表直流电阻替代。

  \item[差分对 Differential Pair]
  用 P/N 两线之间差值传输信息的一组线。它们必须成对布线、控制阻抗、长度和
  参考平面；一条 PCIe lane 包含一对 TX 和一对 RX。

  \item[受控阻抗 Controlled Impedance]
  先依据接口规范确定目标单端/差分阻抗，再由板厂叠层、线宽、间距、铜厚和
  介质参数共同实现。EDA 中写下 ``90 \(\Omega\) differential'' 只是设计
  目标；生产前要让板厂确认可制造参数，必要时用测试 coupon 与 TDR 验证。

  \item[等长 / 时延匹配 Length / Delay Matching]
  让必须同时到达的信号把传播时间差控制在协议预算内。它不是要求全板所有线
  一样长，也不是机械长度越精确越好；应按差分对内、同一字节组或协议 lane
  的规则匹配，并把过孔、连接器和介质传播速度计入。

  \item[AC 耦合 AC Coupling]
  在高速通道串联电容，隔离两端直流共模而让变化信号通过。协议会规定由发送端
  还是接收端提供以及允许数值，不能凭图形对称随意添加。

  \item[连接器 Connector]
  提供可插拔机械、电气边界的器件。接口名称不仅包括数据脚，还包括电源、地、
  屏蔽、检测和控制；能插入不代表引脚定义、电压或协议兼容。

  \item[NC / DNP]
  NC 表示引脚明确不连接；DNP 表示板上可能保留焊盘但默认不装器件。两者都应
  是经过选择的配置，而不是掩盖未完成连线。

  \item[网络类别 Net Class]
  EDA 为一组网络绑定线宽、间距、过孔、差分和长度规则的配置。电源、高速、
  普通信号不能只靠工程师记忆区别；net class 把要求交给 DRC 自动检查，但
  分类错误仍会让正确规则检查错对象。

  \item[BOM]
  Bill of Materials，物料清单。它列出位号、数量、制造商料号、封装、参数、
  可替代料和 DNP 状态；只写“10 k 电阻”通常不足以采购，因为精度、功率、
  封装、温漂和供货状态都可能影响成品。

  \item[Gerber / 钻孔文件]
  板厂通常依据各铜层、阻焊、丝印、板框的 Gerber 文件以及 NC Drill 数据
  制作裸板。这些是某次导出的制造快照，不会随原理图自动更新；下单前必须用
  独立查看器核对层、孔、开窗、板框和版本。

  \item[坐标文件 / 装配图]
  Pick-and-place 坐标记录器件中心、面别、旋转和位号，装配图辅助确认方向与
  特殊工艺。它们要与 BOM、钢网、PCB 和元件 pin 1 标记来自同一版本，否则
  即使裸板正确，也可能把器件旋转或装到错误位置。

  \item[ERC / DRC]
  ERC 检查原理图电气规则，DRC 检查 PCB 几何和设计规则。自动检查能发现大量
  结构错误，但不能替代器件规格、信号完整性、热、机械和制造审查。

  \item[SI / PI]
  Signal Integrity 研究信号在真实互连上是否仍满足电平、时序、眼图和抖动
  预算；Power Integrity 研究电源分配网络的压降、纹波、瞬态和阻抗。连线
  “没有断”只证明拓扑，不能证明高速信号和瞬态供电质量。

  \item[EMI / EMC]
  EMI 是设备产生或受到的电磁干扰，EMC 是设备在规定电磁环境中正常工作且
  不过度干扰他人的能力。回路面积、边沿速度、屏蔽、机壳接地、线缆和滤波
  都会影响结果；通过功能测试不等于通过法规或兼容性测试。
\end{description}

\section*{三、处理器与存储}

\begin{description}[style=nextline,leftmargin=2.8em,labelindent=0em]
  \item[bit / byte]
  bit 是只能表示 0 或 1 的二进制位；byte 是操作系统与多数存储接口使用的
  最小可寻址字节，本书取 1 byte = 8 bits。容量中的 KiB 是 \(1024\) bytes，
  而硬盘厂商常用十进制 kB/GB，读数不同不表示数据凭空丢失。

  \item[十六进制 Hexadecimal]
  用 0--9、A--F 表示基数 16 的数。一个十六进制数字恰好对应 4 bits，所以
  \code{0x3F8}、页表位掩码和机器码用十六进制更易分组；前缀 \code{0x}
  表示记数方式，不是地址本身的特殊类型。

  \item[时钟 Clock]
  周期变化的参考信号，用来协调数字状态机何时采样和推进。频率不是“每秒完成
  多少条指令”的直接等号；PLL、分频、时钟门控和接口参考时钟会形成多个时钟域，
  跨域数据必须同步，参考时钟不稳定时设备即使有电也不能可靠工作。

  \item[CPU]
  Central Processing Unit，中央处理器。它按架构规则取指、译码、执行并提交
  状态变化；现代 CPU 内部会并行和乱序处理，但最终必须呈现规定的软件可见结果。

  \item[SoC]
  System on Chip，把 CPU 核、缓存、内存控制器、PCIe/USB/显示等控制器中的
  多项功能集成到单颗芯片。SoC 并不等于完整电脑：它仍需要供电、时钟、DRAM、
  Flash、连接器和载板，且其中某个控制器存在也不保证板上已经引出相应接口。

  \item[计算模块 / 载板]
  计算模块把 SoC、DRAM、固件存储和部分电源做成可插拔核心；载板把模块触点
  连接到稳压器、连接器和外设。两者共同形成机器，模块规格说明“可能提供什么”，
  载板原理图才说明“本产品怎样实际接出”；SODIMM 外形也不代表 DDR 电气协议。

  \item[架构 Architecture]
  软件能够依赖的指令、寄存器、地址、异常和模式契约。微架构是某款处理器在
  内部怎样实现这些契约；同为 x86-64 的不同 CPU 可有不同流水线和缓存结构。

  \item[寄存器 Register]
  CPU 内部少量、固定宽度且可由指令或控制逻辑直接使用的状态。RIP 指示取指
  位置，RSP 指向栈，CR3 指向页表根；寄存器不等同于大容量 RAM。

  \item[指令 Instruction]
  CPU 架构定义的一次操作，例如移动、加法、跳转或端口访问。汇编文本是可读
  表示，CPU 真正读取的是编码后的机器码字节。

  \item[地址 Address]
  可访问位置的编号，不是位置中的内容。CPU 发出地址、读写方向和宽度，平台
  译码决定请求落到 RAM、ROM、设备还是空洞。

  \item[物理地址 Physical Address]
  经过所有 CPU 地址翻译后送到平台互连的地址。物理地址空间不只包含 RAM，
  还可映射固件、PCI BAR、APIC 和其他设备窗口。

  \item[虚拟地址 Virtual Address]
  分页开启后，指令直接使用并由页表转换的地址。不同进程可在相同虚拟地址看到
  不同物理页，内核也可用高半区映射组织自己的地址空间。

  \item[RAM]
  Random Access Memory，随机存取存储器。运行时可按地址快速读写，常保存
  代码、变量、栈和页表；主存 DRAM 易失且需控制器初始化、刷新和训练。

  \item[DRAM]
  Dynamic RAM，用会泄漏的微小电荷保存状态，因此要刷新。它密度高、适合主存，
  但上电后必须由平台配置内存控制器并完成时序训练。

  \item[内存控制器 / DRAM 训练]
  内存控制器把 CPU 请求转换为 DRAM 命令、地址和时序。上电训练会寻找各数据线
  的可靠采样窗口并配置延时、电压等参数；训练失败时 CPU 可能仍在固件早期运行，
  却没有可用的大容量 RAM，因此看不到普通内核日志。

  \item[SRAM]
  Static RAM，持续供电时无需 DRAM 式刷新。它通常更快但每位成本更高，常用于
  CPU 缓存或芯片内部小存储；“静态”仍不表示断电保留。

  \item[ROM]
  Read-Only Memory。历史上指制造时固定、运行时只读的存储器；现代启动语境
  常泛指被映射为固件入口的非易失区域，物理介质通常是可受控更新的 Flash。

  \item[Flash]
  断电保存、可经擦除和编程协议更新的非易失存储。SPI NOR Flash 常保存平台
  固件；正常软件不能把它当普通 RAM 任意逐字节覆盖。

  \item[缓存 Cache]
  在 CPU 与较慢存储层之间保存近期数据副本的快速存储。缓存提高平均性能，
  也带来一致性和设备 DMA 可见性问题；它不是磁盘文件缓存的唯一含义。

  \item[块设备 Block Device]
  以固定或协议定义的块读写持久数据的设备，例如 SATA/NVMe SSD。CPU 通常
  不能把磁盘扇区直接当普通内存取指，必须经控制器和驱动把数据搬到 RAM。

  \item[SSD]
  Solid-State Drive，以 NAND Flash 保存持久数据，并由控制器完成地址转换、
  纠错、坏块管理、磨损均衡和缓存。SSD 不是“一大块可直接寻址的 RAM”；
  主机通过 SATA 或 NVMe 命令读写逻辑块，断电一致性还取决于 flush 与设备实现。

  \item[扇区 / 逻辑块 / LBA]
  扇区是磁盘历史上的基本传输单位，现代设备向软件报告 Logical Block；
  LBA 是这个逻辑块的序号。例如 512-byte LBA 8 从字节偏移 4096 开始。
  文件系统 block 可由多个设备逻辑块组成，三者不能只因都叫“块”就混用。

  \item[NVMe]
  Non-Volatile Memory Express，面向 PCIe 非易失存储的寄存器和队列协议。
  系统先完成 PCIe 枚举和 BAR 映射，再配置管理/IO 提交与完成队列以及 DMA；
  M.2 插槽有 PCIe 铜线不等于当前内核已经具备 NVMe 驱动。
\end{description}

\section*{四、固件、平台与外设协议}

\begin{description}[style=nextline,leftmargin=2.8em,labelindent=0em]
  \item[固件 Firmware]
  与硬件平台紧密结合、通常保存在 Flash 的最早期软件。它建立处理器、DRAM
  和关键外设条件，再向引导器或操作系统发布内存图与平台信息。

  \item[BIOS / UEFI]
  BIOS 常指传统 PC 固件及其历史接口；UEFI 是现代固件接口和执行环境。二者
  都不是 CPU 硬件本身，也不等同于操作系统内核。

  \item[引导器 Bootloader]
  找到、验证并装入内核，建立执行模式、内存布局和交接信息的软件阶段。本项目
  QEMU 路线的 Stage 1 自行读取 ATA，而实体适配可以新增 UEFI 引导入口。

  \item[总线 / 链路 Bus / Link]
  一组线路及其寻址、时序、编码、仲裁和错误规则。PCIe 是点到点串行链路，
  I2C 是可共享两线总线；“总线”不只是画在图上的粗线。

  \item[HSIO / Pin Mux]
  HSIO 是平台可复用的 High-Speed I/O 通道资源，pin mux 则选择某组物理触点
  在本次启动中承担哪种功能。固件设置、模块能力和 PCB 铜线必须三者一致；
  改 BIOS 复用不会让已经焊好的 USB 铜线自动变成另一连接器上的 PCIe。

  \item[Lane / REFCLK / Link Training]
  高速串行 lane 是一组发送和接收差分对；PCIe 还依赖参考时钟、复位及可能的
  CLKREQ/WAKE 控制。训练期间两端协商速率、宽度和均衡，只有进入 Link Up，
  软件才可能枚举设备；供电正常而训练失败仍会表现为“系统里没有这块卡”。

  \item[设备寄存器 Device Register]
  外设暴露给软件的控制和状态位置。访问可能启动硬件动作、清除状态或产生副作用，
  所以不能把它当普通缓存 RAM 随意重排或重复读取。

  \item[Port I/O]
  x86 独立 I/O 端口空间，通过 IN/OUT 指令访问。COM1 的 \code{0x3F8} 属于
  这种历史接口，与物理内存地址 \code{0x3F8} 不是同一空间。

  \item[MMIO]
  Memory-Mapped I/O，把设备寄存器放入物理地址空间，使用加载/存储访问。
  地址形式像内存，但仍有宽度、顺序、缓存属性和副作用约束。

  \item[PCIe]
  PCI Express，高速点到点分组互连。链路训练以后，软件还要枚举配置空间、
  分配 BAR、设置中断和装入驱动，物理连线成功不等于设备已经可用。

  \item[PCI 配置空间 / BAR]
  每个 PCI/PCIe function 暴露标准配置空间，包含厂商/设备 ID、能力和资源
  请求。BAR 描述设备需要的 MMIO 或 I/O 窗口；固件/操作系统分配地址后，
  驱动才用该窗口访问寄存器，BAR 里的地址不是 PCB 上某根固定“地址线”。

  \item[M.2 M Key / E Key]
  M.2 同时规定卡形、尺寸、缺口位置与可选信号，key 只限制机械/信号组合，
  不是“所有 M.2 都是 NVMe”。本案例 M Key 接 PCIe x1 且不接 SATA；
  E Key 组合卡常让 Wi-Fi 走 PCIe、Bluetooth 走 USB 2.0，两部分需不同驱动。

  \item[USB]
  由主控制器调度的外设总线。连接检测、复位、速度协商、地址、描述符、配置和
  类驱动共同完成枚举；USB 键盘还需要 HID 解析，不能复用 PS/2 寄存器驱动。

  \item[xHCI / HID]
  xHCI 是现代 USB 主控制器接口，操作系统通过 MMIO、环形队列、中断和 DMA
  调度 USB 事务；HID 是键盘、鼠标等人机接口设备类。实体 USB 键盘路径是
  xHCI \(\rightarrow\) USB 枚举 \(\rightarrow\) HID 报告解析，而不是
  QEMU i8042/PS2 端口路径的自然延伸。

  \item[HDMI]
  数字音视频接口。除 TMDS 高速数据外还有 5 V、HPD 和 DDC；源端先发现显示器、
  读取 EDID 并配置模式，之后才发送有效像素。

  \item[EDID / GOP / Framebuffer]
  EDID 是显示器经 DDC 提供的能力数据；UEFI GOP 是固件设置图形模式并交给
  引导软件帧缓冲的接口；framebuffer 是一片按像素格式解释的内存。三者分别
  回答“显示器支持什么”“固件怎样选模式”“软件往哪里写像素”，不能互相代替。

  \item[Ethernet MAC / PHY / Magnetics / RJ45]
  MAC 处理帧、地址与 DMA，PHY 把数字数据转换为线缆上的物理编码，隔离磁性
  器件提供变压器耦合与共模处理，RJ45 是外部连接器。RTL8111H 集成 MAC/PHY
  控制器仍要经过 MDI 差分对和磁性器件才能接网线；驱动还需 PCIe 枚举、
  BAR、描述符环和中断，不能把“网口灯亮”当作网络栈已经工作。

  \item[I2C]
  常用 SCL/SDA 两线连接低速芯片的同步总线。开漏结构依靠上拉回到高电平，
  事务包含起始、地址、应答、数据和停止；I2C 地址不是 CPU 内存地址。

  \item[UART]
  异步串行收发器。双方约定波特率和帧格式，用 TX/RX 发送字节；3.3 V UART
  与使用正负电压的 RS-232 不是可直接相连的电平标准。

  \item[SPI]
  同步串行接口，常包含时钟、片选和一到四根数据线。平台固件 Flash 常用 SPI
  NOR，但能连接一颗 Flash 不表示任意固件镜像都满足平台初始化契约。

  \item[PWM / TACH]
  PWM 通过改变周期中高电平占空比控制风扇目标功率或转速，TACH 是风扇返回的
  转速脉冲。四线风扇的电源、地、PWM 和 TACH 各有职责；软件需知道脉冲数/
  转、极性和频率，不能把 PWM 百分比直接当实际 RPM。

  \item[RTC / 备用电池]
  Real-Time Clock 在主电源关闭时依靠小型备用电源继续计时并保存少量状态。
  它给出墙上时间的起点，不是调度器的高精度单调时钟；电池缺失、晶振停振或
  固件未校准会使日期复位，但不必然阻止 CPU 启动。

  \item[NTC 热敏电阻]
  Negative Temperature Coefficient 电阻随温度升高而减小，常与固定电阻
  构成分压供 ADC 测量。读数要结合具体阻值曲线、上拉电压和 ADC 参考换算；
  它是温度传感元件，不会自己执行风扇控制算法。

  \item[中断 Interrupt]
  设备或定时器异步请求 CPU 处理事件的机制。控制器选择向量和优先级，入口保存
  现场，驱动确认来源并处理；忘记确认可能造成重复中断。

  \item[异常 Exception]
  CPU 执行当前指令时同步检测到的条件，例如非法指令、除零或页故障。异常入口
  与外部 IRQ 可共享 IDT 机制，但原因、返回条件和错误码不同。

  \item[DMA]
  Direct Memory Access，设备直接读写 RAM 的机制。驱动仍要配置地址、长度和
  所有权，并处理缓存一致性、IOMMU、完成通知和失败回收。

  \item[ACPI]
  固件向操作系统描述处理器、中断控制器、电源、定时器和设备拓扑的一组表与
  方法。它不是一个单独芯片，实体平台适配不能用 QEMU 固定地址假设替代。

  \item[APIC]
  Advanced Programmable Interrupt Controller，现代 x86 中断控制体系。
  本地 APIC 属于 CPU 侧，I/O APIC 路由外部中断；它与历史 8259A PIC 有兼容
  关系但不是同一设备。
\end{description}

\section*{五、操作系统与虚拟机器}

\begin{description}[style=nextline,leftmargin=2.8em,labelindent=0em]
  \item[内核 Kernel]
  以高权限长期运行、管理 CPU、内存、设备和隔离的操作系统核心。它不是固件，
  也不是所有用户程序的集合；固件交接后内核逐步接管平台资源。

  \item[驱动 Driver]
  把设备寄存器、队列、中断和错误协议封装为内核服务的软件。驱动必须匹配实际
  控制器；PS/2 驱动不能因为两端都是键盘就自动驱动 USB HID。

  \item[进程 Process]
  一组资源所有权和隔离边界，通常拥有地址空间、文件表与至少一个线程。相同
  虚拟地址在不同进程中可映射不同物理页。

  \item[线程 Thread]
  可被调度的一条执行流，拥有寄存器现场和栈。多个线程可共享一个进程的地址
  空间和文件表，却需要独立保存 CPU/FPU 上下文。

  \item[Ring 0 / Ring 3]
  x86 权限级中的内核与普通用户级。Ring 3 代码不能直接执行特权操作，必须经
  系统调用进入受控内核入口；权限来自处理器检查，不只是编码约定。

  \item[系统调用 System Call]
  用户程序请求内核服务的受控入口。入口必须验证用户指针、长度、对象类型和
  权限，并在返回前恢复正确用户现场；它不是普通跨库函数调用。

  \item[页 Page 与页表 Page Table]
  页是分页管理的固定大小单位，页表把虚拟页映射到物理页并附带读写、用户、
  执行等权限。CR3 指向当前页表根，TLB 缓存近期翻译。

  \item[栈 Stack]
  按后进先出习惯保存调用返回地址、局部状态和中断现场的一段 RAM。CPU 复位后
  不会替裸机固件建立可用栈，早期代码必须先选择并设置 RSP。

  \item[堆 Heap]
  供动态分配器管理的内存区域。它不是磁盘，也不是“所有剩余 RAM”；分配器
  仍需建立范围、对齐、所有权和释放规则。

  \item[文件系统 File System]
  把块设备上的字节组织成超级块、位图、inode、目录和文件等持久结构。缓存、
  写回顺序、校验和和崩溃恢复决定重启后是否仍能解释数据。

  \item[文件描述符 File Descriptor]
  进程文件表中的小整数索引。它引用内核对象或共享打开状态，不是磁盘地址；
  复制描述符、独立打开和关闭会产生不同引用计数与偏移语义。

  \item[QEMU]
  实现目标机器架构和设备行为的模拟器。它可以给 STUOS 提供 x86-64 PC、
  RAM、ROM、串口和 IDE，但不替固件写寄存器，也不等同于 LattePanda 实体
  载板的电气网络。

  \item[TCG]
  QEMU 的 Tiny Code Generator，把客户机指令翻译成宿主可执行代码。宿主即使
  是 AArch64，目标仍按 x86-64 寄存器、机器码和异常规则运行。
\end{description}

\section*{六、六条最值得反复回查的完整链}

\begin{enumerate}[leftmargin=2em]
  \item \textbf{冷启动}：输入电源 \(\rightarrow\) 稳压/时钟
        \(\rightarrow\) 复位释放 \(\rightarrow\) ROM/Flash 取指
        \(\rightarrow\) DRAM 初始化 \(\rightarrow\) 引导器
        \(\rightarrow\) 内核。
  \item \textbf{程序装入}：磁盘块 \(\rightarrow\) 控制器/驱动
        \(\rightarrow\) RAM \(\rightarrow\) ELF 校验
        \(\rightarrow\) 页表 \(\rightarrow\) 用户 RIP/RSP。
  \item \textbf{键盘输入}：按键 \(\rightarrow\) 物理接口
        \(\rightarrow\) 控制器队列 \(\rightarrow\) 中断
        \(\rightarrow\) 驱动 \(\rightarrow\) FIFO
        \(\rightarrow\) fd 0 \(\rightarrow\) Shell。
  \item \textbf{地址访问}：虚拟地址 \(\rightarrow\) 页表/TLB
        \(\rightarrow\) 物理地址 \(\rightarrow\) 平台译码
        \(\rightarrow\) RAM、ROM、MMIO 或故障。
  \item \textbf{实体 PCIe 设备}：差分铜线/时钟/复位
        \(\rightarrow\) 链路训练 \(\rightarrow\) 配置空间
        \(\rightarrow\) BAR/中断 \(\rightarrow\) 驱动
        \(\rightarrow\) 内核服务。
  \item \textbf{错误证据}：条件未建立 \(\rightarrow\) 首个断点
        \(\rightarrow\) 电压/寄存器/日志/退出码
        \(\rightarrow\) 分层定位，而不是把所有失败都叫“死机”。
\end{enumerate}
